\documentclass[11pt]{article}
\usepackage[margin=1in]{geometry}
\usepackage{booktabs}
\usepackage{amsmath}
\usepackage{hyperref}
\usepackage{graphicx}
\usepackage{microtype}
\hypersetup{colorlinks=true, linkcolor=blue!60!black, citecolor=blue!60!black, urlcolor=blue!60!black}

\title{Neural Weight Compression: Lossless BF16 Weights Decoded Inside the Matrix--Vector Kernel}
\author{Paul Otto\\ cloo GmbH\\ \texttt{https://github.com/parda21/NWC}}
\date{September 2026 --- draft}

\begin{document}
\maketitle

\begin{abstract}
Token generation with large language models is bound by memory bandwidth: every weight is read once per
token. We present NWC, a lossless compression scheme for BF16 weights that is decoded inside the
matrix--vector kernel, in registers, so that no decompressed weight is ever written to memory. The
mantissa byte of a BF16 weight carries 7.97 bits of entropy and is stored raw; the exponent/sign byte
carries 2.6--2.9 bits and is coded with a frequency-ranked prefix code that a 4096-entry shared-memory
lookup table decodes two symbols at a time. The compressed model is 31\,\% smaller and bit-exact.
Because the fused kernel reads only the compressed bytes, it is faster than the uncompressed cuBLAS
matvec on GPUs whose decoder throughput exceeds \emph{bandwidth}$\times$0.69: on an RTX~4070, Qwen3-4B
generates 55 instead of 45 tokens/s at 5.67 instead of 8.10\,GB of VRAM with bit-identical logits; on a
bandwidth-starved NVIDIA A16 (10 SMs) it generates 18.2 instead of 16.8 tokens/s. DFloat11, which
compresses the same bits but decodes into a buffer before the matmul, takes 2.7$\times$ the native GPU
time per token on the same hardware. We report the instruction-level cost model that made the in-kernel
decoder fast enough on Ampere, and the negative results along the way.
\end{abstract}

\section{Introduction}

Serving a transformer at batch size one is a sequence of matrix--vector products in which each weight is
read exactly once. Time per token is therefore set by the number of bytes the memory system delivers, not
by arithmetic. Quantization attacks the byte count by changing the weights; lossless compression attacks
it without changing them, but only pays off if decompression does not itself become the bottleneck.
Existing lossless approaches for LLM weights \cite{dfloat11,zipnn} decompress into a buffer that the
subsequent matmul reads again, so the decompressed weights cross the memory system twice and batch-one
latency gets worse, not better.

NWC (Neural Weight Compression) takes the other route: the weights stay compressed in VRAM and are
decoded inside the matvec kernel, in the registers of the thread that immediately multiplies them. The
kernel reads 69\,\% of the bytes and never writes a decompressed weight. Whether this is a net win is a
question of decoder throughput per streaming multiprocessor versus memory bandwidth; the contribution of
this paper is a format and a decoder that win on consumer GDDR parts by a wide margin and, after a
sequence of measured redesigns, also on an Ampere part with 10 SMs.

Contributions: (1) a byte-plane format with a rank prefix code for the exponent byte whose compressed size
matches the entropy limit within 1--2\,\%; (2) a fused CUDA decoder that needs 11--13 SM clocks per 64
weights on Ampere, including a lookup-table design that yields two codes per shared-memory access and a
bit-window refill without a counter; (3) an instruction-level cost model of the Ampere GA10x SM, measured
rather than assumed, which predicted the wins and explains the losses; (4) end-to-end measurements on two
GPUs with bit-exact checks, and a head-to-head comparison with DFloat11.

\section{Related work}

DFloat11 \cite{dfloat11} observes that the BF16 exponent is highly redundant and Huffman-codes it,
reaching 70\,\% of the original size with bit-exact reconstruction; decoding runs in a separate kernel into
a BF16 buffer, with hierarchical lookup tables in shared memory. ZipNN \cite{zipnn} applies byte-grouping
and general-purpose entropy coding to checkpoints for storage and transfer. Both agree with our entropy
measurement: the mantissa is incompressible, the exponent carries roughly three bits. Asymmetric numeral
systems \cite{ans} were our first coder (formats v4--v8, Section~\ref{sec:history}); the final format uses a
prefix code because its decoder needs no table gather per symbol. Weight quantization (INT4/INT8, FP8) is
orthogonal: it changes the weights and trades quality for size, whereas NWC is lossless; the two can be
combined only to the extent that quantized weights still carry redundancy, which they largely do not.

\section{Method}

\subsection{What is compressible}

A BF16 value has one sign bit, eight exponent bits and seven mantissa bits. Over the linear layers of
Qwen2.5 and Qwen3 models we measure 7.97 bits of entropy in the mantissa byte and 2.6--2.9 bits in the
byte formed by sign and the upper seven exponent bits. Context modelling (neighbouring weights, row
statistics, strided correlation) recovers less than 0.4\,\% and costs more in table traffic than it saves.
The achievable size is therefore $(8 + 2.8)/16 \approx 0.68$ of BF16, and this bound is the same for every
entropy coder; the design freedom is entirely in decoding cost.

\subsection{Format}

The exponent byte is mapped to its rank in the per-tensor frequency order. Rank $r < 15$ is written as
$r$ zeros, a one, and the raw sign bit ($r+2$ bits); rank $\geq 15$ is an escape of 16 bits followed by
the raw byte. The mean cost is 2.85--2.99 bits per weight on Qwen models (escapes 0.01--0.4\,\%); the rANS
coder of earlier formats saved about 1\,\% of the exponent stream at roughly twice the decoding cost.

The matrix is tiled into blocks of $8$ rows $\times$ $512$ columns. One warp decodes one block; lane $L$
owns columns $16L \ldots 16L+15$ of all eight rows and has its own bit stream. Streams are packed into
32-bit words, 32 streams per block, with a one-byte length per lane and a warp prefix sum for the offsets.
Matrices of any shape are supported: the block grid is padded to multiples of 8 and 512, padding weights
cost two bits and have no mantissa, and the mantissa plane is stored row-major with the row length
rounded up to 16.

\subsection{Decoder}

\paragraph{Lookup table.} A table of 4096 32-bit entries, indexed by the next 12 bits of a lane's stream,
holds for 98.9\,\% of the patterns two complete codes: the two exponent bytes and the total number of bits
consumed, placed in the low byte so that the funnel-shift instruction can use the entry directly as its
shift amount. Patterns without two complete codes are flagged and finished by a \texttt{clz}-based slow
path (0.06--0.4\,\% of lane steps). The table lives in shared memory; the rank-to-exponent mapping is
recovered from the table itself, so a block needs a single 16\,KB copy.

\paragraph{Bit window.} Each lane keeps a 64-bit window. Advancing by $c$ bits is two funnel shifts. The
fill level is tracked as a rotated power of two, $M = 2^{64-n_b} \bmod 2^{32}$, rotated left by $c$ each
step; a wrap-around signals that fewer than 32 valid bits remain, and the pre-loaded next word $w$ is
merged with a single wide multiply-add, $\mathit{hi{:}lo} \mathrel{+}= w \cdot M$, which is a no-op for
lanes that did not wrap. Only the load of the following word is predicated.

\paragraph{Weights and accumulation.} One \texttt{prmt} assembles the fp32 bit pattern of a weight from its
exponent byte and its mantissa byte, using sign replication of a known-zero byte for the two zero bytes.
The 16 activation values of a lane are loaded once per block into registers; each row is accumulated in
a register and a transposed shuffle reduction at block end leaves the sum of row $r$ in lane $4r$. The
kernel runs one persistent block of 1024 threads per SM. Dequantization and the embedding lookup for the
tied output projection use the same decoder.

\subsection{Cost model}
\label{sec:cost}

Nsight Compute was not available on the measurement machines, so the decoder was optimised against
knock-out variants and a micro-benchmark of per-instruction issue cost. On the Ampere GA10x SM, with
independent instruction chains and data-dependent operands (with constant operands \texttt{ptxas} folds
the chains and reports twice the rate), we measure in SM clocks per warp instruction: FFMA 0.27;
IADD3/LOP3 0.30; IMAD, IMAD.WIDE, SHF, ISETP, SEL, PRMT 0.51; IMAD.HI 1.0; FLO 2.0. IMAD and SHF overlap
(0.26 when interleaved), as do SHF and IADD3; PRMT and IMAD.HI block both pipes. The issue limit of one
instruction per clock per sub-partition bounds a 21-instruction step at 5.5 clocks; the measured 11--13
clocks per 64 weights are consistent with the pipe costs plus the branch and memory instructions.

Two effects dominated the path from 17.7 to 11.0 clocks (Table~\ref{tab:stages}): code size (a fully
unrolled kernel with the slow path inline was 147\,KB and cost 3.5 clocks per step against a 23\,KB rolled
loop), and moving the window arithmetic from the IMAD pipe to the shift pipe once the true issue rates
were known.

\section{Experiments}

\subsection{Setup}

RTX~4070 (12\,GB GDDR6X, 452\,GB/s measured, 46 SMs, CUDA 13.4) and an NVIDIA A16 in a vGPU 16Q profile
(16\,GB GDDR6, 10 SMs, 2\,MB L2, cuBLAS reaches 160--165\,GB/s; CUDA 12.6, sm\_86). Model: Qwen3-4B in BF16
via HF Transformers, with q/k/v and gate/up projections fused and the tied \texttt{lm\_head} compressed.
Kernel times are medians of CUDA-event measurements with weight copies rotated against the L2; token
rates are wall-clock over 256 greedy tokens with the decode step captured as a CUDA graph. Correctness is
checked by bit-exact dequantization and gather tests for a range of shapes including $K = 11008, 3600,
1000, 4097$ and $M = 1001, 13, 3$, and by comparing the generated token sequence with HF
\texttt{generate}.

\subsection{Results}

\begin{table}[h]\centering
\begin{tabular}{lll}\toprule
 & RTX 4070 & NVIDIA A16 \\\midrule
VRAM, native $\rightarrow$ NWC & 8.10 $\rightarrow$ 5.67\,GB & 8.10 $\rightarrow$ 5.67\,GB \\
tokens/s, native $\rightarrow$ NWC & 45.0 $\rightarrow$ 55.2 & 16.8 $\rightarrow$ 18.2 \\
GPU time per token, native $\rightarrow$ NWC & 20.4 $\rightarrow$ 18.9\,ms & 54.1 $\rightarrow$ 52.1\,ms \\
weight kernels vs cuBLAS (5 shapes) & 1.17--1.55$\times$ & 1.06--1.18$\times$ \\
logits vs native & bit-identical & max 0.20, argmax identical$^{a}$ \\
token sequence vs HF generate (64) & identical & identical \\\bottomrule
\end{tabular}
\caption{Qwen3-4B, greedy decoding as a CUDA graph. $^{a}$cuBLAS itself differs by max 0.19 between the
two GPUs (different fp32 summation order); the weights are bit-exact on both.}
\end{table}

\begin{table}[h]\centering
\begin{tabular}{lrrrr}\toprule
layer ($M \times K$) & 4070 cuBLAS / NWC & ratio & A16 cuBLAS / NWC & ratio \\\midrule
qkv $6144 \times 2560$ & 0.117 / 0.100\,ms & 1.17 & 0.200 / 0.178\,ms & 1.12 \\
o $2560 \times 4096$ & 0.086 / 0.073\,ms & 1.18 & 0.135 / 0.128\,ms & 1.06 \\
gate\_up $19456 \times 2560$ & 0.252 / 0.163\,ms & 1.55 & 0.605 / 0.513\,ms & 1.18 \\
down $2560 \times 9728$ & 0.146 / 0.110\,ms & 1.34 & 0.306 / 0.271\,ms & 1.13 \\
lm\_head $151936 \times 2560$ & 1.675 / 1.241\,ms & 1.35 & 4.787 / 4.538\,ms & 1.06 \\
sum per token & 23.4 / 17.3\,ms & 1.35 & 49.6 / 43.8\,ms & 1.13 \\\bottomrule
\end{tabular}
\caption{Per-layer matvec time versus cuBLAS.}
\end{table}

\subsection{Comparison with DFloat11}

Same model (Qwen/Qwen3-4B and DFloat11/Qwen3-4B-DF11, dfloat11 0.5.0), same prompt, same profiler, RTX
4070: native 20.4\,ms GPU time per token, DFloat11 56.1\,ms (2.74$\times$; its decode kernel alone takes
36.6\,ms), NWC 18.9\,ms. VRAM 8.10 / 5.73 / 5.67\,GB. The size is the same because it is the entropy of the
mantissa; the speed differs because DFloat11's decoded weights pass through memory twice.

\subsection{Decoder stages and knockouts}
\label{sec:history}

\begin{table}[h]\centering
\begin{tabular}{lrrl}\toprule
stage & gate\_up & lm\_head & change \\\midrule
v8 (rANS on exponent pairs, 64-bit state) & 17.7 & 18.3 & --- \\
prefix code, 12-bit LUT, IMAD window & 17.5 & 19.8 & no table gather, INT pipe busy \\
8$\times$512 block, activations in registers & 16.0 & 17.2 & no activation loads \\
slow path out of line & 12.5 & 14.3 & code 147\,KB $\rightarrow$ 52\,KB \\
row loop not unrolled & 12.5 & 14.0 & 23\,KB, slow path inline \\
funnel-shift window, rotation fill level & 12.1 & 13.3 & work moved to the shift pipe \\
L1 prefetch of the block stream & \textbf{11.0} & \textbf{12.6} & \\\bottomrule
\end{tabular}
\caption{SM clocks per warp step of 64 weights on the A16; parity with cuBLAS is 13.6.}
\label{tab:stages}
\end{table}

Knockouts of the final kernel (gate\_up, 12.1 without prefetch): without the slow path 11.0, without
refills 9.7, without the table access 10.9, floor 9.7. Variants that did not help: a warp-uniform branch
via \texttt{\_\_any\_sync} ($+1$ clock), unrolling the row loop ($\pm 0$ / $+0.3$), 13/14-bit peeks
(slower), IMAD.HI instead of SHF ($\pm 0$). Earlier, with rANS: 64-bit table entries ($+4\,\%$), four states
per lane ($-5$ to $-15\,\%$), 8192-weight blocks (always slower: the streams of 32--40 warps thrash the
L1), quad symbols (exponent bytes are nearly independent: 2.815 / 2.811 / 2.803 bits per weight for 1/2/4
byte symbols).

\subsection{Where it wins}

The decoder delivers 7--8\,GB/s of compressed data per SM$\cdot$GHz (v8: 5.0, v6: 3.5--3.9). NWC is faster
than native when this exceeds bandwidth $\times$ 0.69. For the 4070 and the 4080 Super this is satisfied by
a factor of three (they are bandwidth-bound); for the A16 by 10--20\,\%; projected for an H100 PCIe
(2\,TB/s) at 1.1--1.2$\times$ and for an H100 SXM (3.35\,TB/s) at 0.8--0.9$\times$, assuming Ampere issue
rates. Closing the SXM gap requires 20\,\% fewer clocks per step; a tensor-core accumulation path in which
one \texttt{prmt} produces the bf16x2 A-fragment of an \texttt{mma.m16n8k16} directly is the planned next
step.

\section{Limitations}

Batch-one decoding is the fast path; prefill dequantizes into a temporary BF16 buffer and only saves
memory. The gain is bounded by the byte ratio (at most 1/0.69) and realised only where the decoder keeps
up with the memory system. Logits are bit-identical to the native computation only where the fp32
summation order matches; otherwise they lie within the spread cuBLAS itself shows between GPUs.
Measurements cover two GPUs and one model family; Hopper has not been measured.

\section{Conclusion}

Lossless weight compression does not have to cost latency. Decoding inside the matvec turns the 31\,\%
size reduction into a speed-up on memory-bound hardware, provided the decoder is designed against the
real instruction costs of the SM. The format, kernel, tests and every measurement in this paper are
available under the Apache-2.0 license.

\section*{Acknowledgements}
This work was sponsored by cloo GmbH.

\begin{thebibliography}{9}
\bibitem{dfloat11} T.~Zhang, Y.~Sui, S.~Zhong, V.~Chaudhary, X.~Hu, A.~Shrivastava. 70\% Size, 100\% Accuracy:
Lossless LLM Compression for Efficient GPU Inference via Dynamic-Length Float. arXiv:2504.11651, 2025.
\bibitem{zipnn} M.~Hershcovitch et al. ZipNN: Lossless Compression for AI Models. arXiv:2411.05239, 2024.
\bibitem{ans} J.~Duda. Asymmetric numeral systems: entropy coding combining speed of Huffman coding with
compression rate of arithmetic coding. arXiv:1311.2540, 2013.
\bibitem{qwen3} Qwen Team. Qwen3 Technical Report. arXiv:2505.09388, 2025.
\end{thebibliography}

\end{document}
